\section{Discussion and Limitations}
OrnaScout is designed for a regime in which exhaustive high-resolution VLM inference is impractical and silent errors are costly. Its potential contribution is therefore not another plant classifier, but a measurable evidence-allocation and abstention policy. The decisive tests are whether the selector improves task performance at equal visual budget, whether controlled exploration contributes beyond pure detector ranking, whether the gate yields useful coverage at a predeclared risk, and whether these effects persist across crop direction and VLM backbone.

The validation evidence audit already narrows the claim. Detector-guided ranking retains more annotated instances than random selection, but the exploit--explore policy does not maximize evidence recall at $B=8$; count top-$K$ does. Exploration instead trades some local recall for lower overlap and greater hard-negative exposure. Whether that trade improves semantic accuracy or selective reliability remains an empirical question for the frozen VLM runs, rather than an assumption built into the method narrative.

Several limitations are material. First, the dataset is imbalanced across crop, source, and annotation density. Second, one-class boxes support localization and count-derived targets but do not provide disease, quality, or management ground truth. Third, the selector inherits detector errors; exploration reduces but cannot remove that dependence. Fourth, the gate is empirically calibrated and may fail under severe target shift. Fifth, peak resident memory on a development laptop is not equivalent to energy or latency on a production edge accelerator. Finally, a private corpus can support controlled evaluation but should not be presented as a public dataset contribution unless release permission and documentation are secured.

\section{Conclusion}
We presented a framework for allocating compact-VLM computation to detector-supported regions and declining uncertain predictions under cross-crop aerial shift. The evaluation is structured around equal budgets, source-safe splits, clustered uncertainty, hard negatives, and measured edge costs. The final scientific conclusion will be determined by the frozen experiments: \tbd{replace this sentence with the supported routing, transfer, and selective-risk findings}.
